\section{Read-Time Defenses for Derived State}
\label{sec:defense}

All defenses below share the same retrieval and response-shape settings: a live IAM check on every
retrieved chunk, ACL filtering inside the vector query, and constant-shape responses. They differ in
how they decide whether a cache entry or a memory item may still be served.

\subsection{Provenance-tainted revocation}
\paragraph{Taint.}
Every derived artifact $a$ carries a taint $T(a)\subseteq\mathcal{D}$, the set of source documents
whose content it may contain. A cache entry has the taint of the answer it stores. An answer has the
union of the taints of the chunks and memory items in its prompt; a summary has the union of the
taints of the turns it replaces:
\begin{equation}
T(\mathrm{answer}) = \bigcup_{c\in\mathrm{context}} T(c), \qquad
T(\mathrm{summary}) = \bigcup_{t\in\mathrm{turns}} T(t).
\end{equation}
The first equation makes taint transitive through memory: an answer built from an earlier answer
about document $d$ inherits $d$.

\paragraph{Enforcement.}
\emph{Eagerly}, when a revocation event fires, tainted cache entries are evicted and tainted memory
items are purged. \emph{Lazily}, on every read, artifact $a$ is served to user $u$ at time $t$ only if
\begin{equation}
\forall d\in T(a):\; \mathrm{IAM}(u,d,t)=\text{allow}.
\label{eq:lazy}
\end{equation}
The lazy check refers to the identity system at read time, so it does not depend on whether a
revocation event reached the artifact. The cache evaluates it only for candidates above the
similarity threshold, from the most similar down, so authorization work grows with the number of
servable candidates rather than with the cache size.

\paragraph{Sentence-level taint.}
A summary's taint is the union of everything it condenses, so one revoked source discards the whole
summary, including material from documents the user still holds. With an extractive summarizer, each
summary sentence can carry the taint of the turn it came from, and the lazy check then drops only the
tainted sentences (\texttt{full\_fine}). This stays sound with any generator, since a filtered
summary's taint is the union of the sentences it keeps. Two things limit it. With model-written
summaries there is no sentence-to-source mapping, so \texttt{full\_fine} reduces to \texttt{full};
this paper does not measure it. And an answer's taint covers every
memory item in its prompt unless the generator reports which items it used: our stub does, real
models do not, so with real models an answer given after the user discussed the revoked document
carries that document's taint even if it did not use it. Context
attribution~\cite{cohenwang2024contextcite} or permissive information-flow
analysis~\cite{siddiqui2024permissive} is the route to finer answer taint, at the price of a
soundness argument we do not make.

\subsection{Simple alternatives}
A reviewer's first question is why taint is needed at all. We implement the two obvious
alternatives with identical retrieval and response settings.
\begin{itemize}
\item \emph{Permission epoch} (\texttt{epoch}). Each user has an epoch in the IAM that every
revocation increments; a reclassification increments the epoch of every user who loses the
document. Cache entries and memory turns carry the epoch of the user who created them
and are dropped on read when it is stale; a summary written after a revocation does not absorb stale
turns. This is a read-time check, so it covers lost events, but it discards everything the user
derived before any revocation, whatever its sources.
\item \emph{Session flush} (\texttt{flush}). On a revocation event, every session and cache entry of
that user is dropped. It is eager and depends on event delivery: a revocation hook that is never
delivered, as happens with webhooks or a message bus, leaves the user's state in place.
\end{itemize}
The read-time alternatives reach zero leakage by construction for the artifacts they cover. The
comparison that matters is therefore what each one costs (\cref{sec:rq5}). The cost differs most
when a document is reclassified: the epoch of every user who loses it is bumped and all their
sessions are discarded, whereas taint discards only what derives from that document. The existence
experiment also runs \texttt{full} with a per-user cache (\texttt{full\_user}), so that the comparison
does not mix the defense with cache scope.

\subsection{Constant-shape responses}
Responses expose relevance scores only in coarse bands and only over chunks the user may see, use one
refusal text for every reason, and pad latency up to the next multiple of a bucket (500\,ms). Padding
is applied after all computation.

\subsection{Security argument and assumptions}
\label{sec:soundness}
\begin{invariant}
For every derived artifact $a$ (cache entry, memory turn, summary or summary sentence), if content
of a restricted document $d$ reached $a$ through retrieval or memory, then $d\in T(a)$.
\end{invariant}

\begin{proposition}
Assume A1--A4 below. If the invariant holds and every read of a derived artifact applies
\eqref{eq:lazy}, then after $u$ loses access to $d$ no response to $u$ is generated from, or
re-serves, an artifact containing content of $d$, and no retrieval for $u$ returns a chunk of $d$.
\end{proposition}
\begin{proof}[Proof sketch]
Content enters an artifact only through a generation step, whose inputs are the retrieved chunks,
the memory items in the prompt and the query. The pipeline writes the union of the chunks' documents
and the items' taints as the response taint, cache and memory writes store it unchanged, a cache hit
re-serves an entry with its stored taint, and a summary's taint (or each summary sentence's) is the
union of the taints it condenses. By induction over the sequence of writes, the invariant holds for
every artifact. After revocation $\mathrm{IAM}(u,d,t)=\text{deny}$ for every read (A4), so any
artifact with $d\in T(a)$ fails \eqref{eq:lazy}, and the live check removes $d$'s chunks from
retrieval. The response is then generated only from artifacts and chunks free of $d$'s content (A1--A3).
\end{proof}

A randomized test runs sessions with frequent summarization, at turn and sentence granularity, and
checks that no artifact holds a secret unique to a restricted document without that document in its
taint; it also confirms that dropping transitivity violates the invariant.

\paragraph{Assumptions.}
\begin{itemize}
\item[A1] \emph{User input is untainted.} Pasted restricted content is not attributed to its source.
An optional switch fingerprints queries against restricted secrets and taints a query containing one
verbatim; paraphrased re-entry is the declassification limit of information-flow
control~\cite{myers1997decentralized,sabelfeld2003language}. Measuring this channel is future work
(\cref{sec:conclusion}).
\item[A2] \emph{No untracked inputs.} Tools, web search or plug-ins would have to label their output.
\item[A3] \emph{Document ACLs are correct.} A document that copies restricted figures under a broader
ACL legitimately serves them.
\item[A4] \emph{Consistent authorization reads.} Checks read IAM state at least as new as the
revocation; a lagging replica reopens the ``new enemy'' window~\cite{pang2019zanzibar}.
\end{itemize}

\paragraph{Presets.}
\Cref{tab:presets} defines the deployments. \texttt{strawman} is a deliberately leaky reference.
\texttt{baseline} follows the documented defaults of \cref{tab:grounding}. \texttt{baseline\_shared}
is the same with a tenant-shared cache, a misconfiguration we measure separately. \texttt{legacy} purges
memory by document identifier, which cannot match a summary. \texttt{authz\_prefilter} and
\texttt{authz\_postfilter} implement the enforcement patterns of permission-aware
RAG~\cite{jeong2025permissionrag} and Zanzibar-style check-after-retrieve~\cite{pang2019zanzibar} on our
switches; they are reimplementations of the pattern, not those systems.

\begin{table*}[t]
\centering
\caption{Deployment presets. ``Stale ACL'' means the index ACL copy is trusted between syncs. All
presets except \texttt{strawman} retrieve within the user's tenant and expose scores computed over
chunks the user may see.}
\label{tab:presets}
\footnotesize
\begin{tabular}{@{}p{2.2cm}p{3.1cm}p{3.6cm}p{3.6cm}p{3.6cm}@{}}
\toprule
Preset & Retrieval & Cache & Memory & Response metadata \\
\midrule
strawman & all tenants, post-filter, withheld scores exposed & global, no hook & no purge & raw scores, distinct refusals \\
baseline & pre-filter, stale ACL & per user, no hook & no purge & raw \\
baseline\_shared & pre-filter, stale ACL & tenant-shared, no check & no purge & raw \\
legacy & pre-filter, stale ACL & tenant, evict on revoke & purge by document ID & bands, uniform refusal, 30\,ms floor, fast refusal \\
authz\_prefilter & pre-filter, live IAM & per user, no hook & no purge & raw \\
authz\_postfilter & post-filter, live IAM & per user, no hook & no purge & raw \\
epoch & pre-filter, live IAM & per user, epoch check & epoch check & bands, uniform refusal, 500\,ms buckets \\
flush & pre-filter, live IAM & per user, flush on revoke & flush on revoke & as epoch \\
full & pre-filter, live IAM & tenant, evict + lazy taint check & transitive taint purge + lazy check & as epoch \\
full\_user & as full & per user, evict + lazy taint check & as full & as epoch \\
full\_fine & as full & as full & as full, sentence-level summary taint & as epoch \\
\bottomrule
\end{tabular}
\end{table*}
